\textbf{Phần 1. Trắc nghiệm lựa chọn}
\setcounterpageref{ex}{0}

\begin{ex} %[3D2V25]
Trong một phép chia có số chia là 7, số dư lớn nhất có thể có của phép chia đó là
\choice
{\True $6$}
{$7$}
{$5$}
{$4$}
\loigiai{
Số dư luôn luôn bé hơn số chia. Với số chia là 7 thì số dư lớn nhất là $7 - 1 = 6$.
\textbf{Đáp án của câu hỏi là A}
}
\end{ex}

\begin{ex} %[3D2V25]
Thực hiện phép chia: $47 : 6$. Thương và số dư lần lượt là
\choice
{$8$ và dư $1$}
{$7$ và dư $6$}
{\True $7$ và dư $5$}
{$6$ và dư $5$}
\loigiai{
Ta có: $47 = 6 \times 7 + 5$ (với $5 < 6$). Thương bằng 7 và dư 5.
\textbf{Đáp án của câu hỏi là C}
}
\end{ex}

\begin{ex} %[3D2V25]
Một người có 35 quả cam chia đều vào các túi, mỗi túi có 4 quả. Hỏi cần ít nhất bao nhiêu túi để đựng hết số quả cam đó?
\choice
{$8$ túi}
{\True $9$ túi}
{$10$ túi}
{$7$ túi}
\loigiai{
Ta có: $35 : 4 = 8$ (dư 3 quả). Cần thêm 1 túi nữa để đựng 3 quả còn dư: $8 + 1 = 9$ túi.
\textbf{Đáp án của câu hỏi là B}
}
\end{ex}

\begin{ex} %[3D2V25]
Trong phép chia cho 5, thương bằng 9 và số dư là số dư lớn nhất có thể. Số bị chia của phép chia đó là
\choice
{$45$}
{$48$}
{$44$}
{\True $49$}
\loigiai{
Số dư lớn nhất là 4. Số bị chia là: $5 \times 9 + 4 = 49$.
\textbf{Đáp án của câu hỏi là D}
}
\end{ex}

\begin{ex} %[3D2V25]
Có 45 học sinh đi đò qua sông. Mỗi chuyến đò chỉ chở được nhiều nhất 6 học sinh (không kể người lái đò). Hỏi cần ít nhất bao nhiêu chuyến đò để chở hết số học sinh đó qua sông?
\choice
{\True $8$ chuyến đò}
{$7$ chuyến đò}
{$9$ chuyến đò}
{$6$ chuyến đò}
\loigiai{
Ta có: $45 : 6 = 7$ (dư 3 học sinh). 3 học sinh còn lại cần thêm 1 chuyến đò: $7 + 1 = 8$ chuyến.
\textbf{Đáp án của câu hỏi là A}
}
\end{ex}

\begin{ex} %[3D2V25]
Một đoàn du khách gồm 26 người muốn đi taxi. Mỗi xe taxi chở được 4 người. Cần ít nhất bao nhiêu xe taxi để chở hết đoàn khách?
\choice
{$6$ xe taxi}
{$8$ xe taxi}
{\True $7$ xe taxi}
{$5$ xe taxi}
\loigiai{
Ta có: $26 : 4 = 6$ (dư 2 người). Cần thêm 1 xe nữa: $6 + 1 = 7$ xe.
\textbf{Đáp án của câu hỏi là C}
}
\end{ex}

\begin{ex} %[3D2V25]
Trong một phép chia có dư, số bị chia là 59, số chia là 8. Số dư của phép chia đó là
\choice
{$4$}
{\True $3$}
{$5$}
{$2$}
\loigiai{
Ta có: $59 : 8 = 7$ (dư 3) vì $8 \times 7 + 3 = 59$.
\textbf{Đáp án của câu hỏi là B}
}
\end{ex}

\begin{ex} %[3D2V25]
Bác nông dân có 50 quả bưởi đóng vào các thùng carton, mỗi thùng chứa 6 quả bưởi. Sau khi đóng đầy các thùng, bác còn thừa bao nhiêu quả bưởi?
\choice
{$4$ quả bưởi}
{$3$ quả bưởi}
{$1$ quả bưởi}
{\True $2$ quả bưởi}
\loigiai{
Ta có: $50 : 6 = 8$ (dư 2). Bác đóng được 8 thùng và thừa 2 quả bưởi.
\textbf{Đáp án của câu hỏi là D}
}
\end{ex}

\begin{ex} %[3D2V25]
Một sợi dây dài $75\text{ cm}$ được cắt thành các đoạn dài $8\text{ cm}$. Người ta cắt được nhiều nhất bao nhiêu đoạn như thế?
\choice
{\True $9$ đoạn}
{$8$ đoạn}
{$10$ đoạn}
{$7$ đoạn}
\loigiai{
Ta có: $75 : 8 = 9$ (dư $3\text{ cm}$). Cắt được nhiều nhất 9 đoạn.
\textbf{Đáp án của câu hỏi là A}
}
\end{ex}

\begin{ex} %[3D2V25]
Nếu chia một số tự nhiên cho 4 thì các số dư có thể nhận được là
\choice
{$1, 2, 3, 4$}
{$0, 1, 2$}
{\True $0, 1, 2, 3$}
{$1, 2, 3$}
\loigiai{
Phép chia cho 4 có thể chia hết (dư 0) hoặc dư 1, 2, 3.
\textbf{Đáp án của câu hỏi là C}
}
\end{ex}

\begin{ex} %[3D2V25]
Tìm số nhỏ nhất mà khi chia cho 7 thì có số dư là 5 và thương bằng 6. Số đó là
\choice
{$42$}
{\True $47$}
{$49$}
{$45$}
\loigiai{
Số đó là: $7 \times 6 + 5 = 47$.
\textbf{Đáp án của câu hỏi là B}
}
\end{ex}

\begin{ex} %[3D2V25]
Lớp 3C có 34 bạn xếp thành các hàng, mỗi hàng có 5 bạn. Hỏi có bao nhiêu hàng có đủ 5 bạn và hàng cuối cùng có bao nhiêu bạn?
\choice
{$6$ hàng đủ 5 bạn và hàng cuối có $5$ bạn}
{$7$ hàng đủ 5 bạn và hàng cuối có $1$ bạn}
{$5$ hàng đủ 5 bạn và hàng cuối có $9$ bạn}
{\True $6$ hàng đủ 5 bạn và hàng cuối có $4$ bạn}
\loigiai{
Ta có: $34 : 5 = 6$ (dư 4). Vậy có 6 hàng đủ 5 bạn và hàng cuối cùng có 4 bạn.
\textbf{Đáp án của câu hỏi là D}
}
\end{ex}

\textbf{Phần 2. Trắc nghiệm đúng sai}
\setcounterpageref{ex}{0}

\begin{ex} %[3D2V25]
Xét tính đúng sai của các khẳng định về phép chia có dư:
\choiceTF
{\True Trong mọi phép chia có dư, số dư luôn luôn bé hơn số chia.}
{Trong phép chia cho 6, số dư lớn nhất là 6.}
{\True Số bị chia bằng thương nhân với số chia rồi cộng với số dư.}
{Một số chia cho 5 dư 3 thì cộng thêm 1 đơn vị sẽ chia hết cho 5.}
\loigiai{
\begin{itemize}
    \item Ý a) Đúng: Định lí phép chia có dư.
    \item Ý b) Sai: Sai vì số dư lớn nhất là $6 - 1 = 5$.
    \item Ý c) Đúng: Công thức $a = b \times q + r$.
    \item Ý d) Sai: Sai vì dư 3 cộng thêm 2 đơn vị mới chia hết cho 5.
\end{itemize}
\begin{itemize}
    \item a) \textbf{Đúng}.
    \item b) \textbf{Sai}.
    \item c) \textbf{Đúng}.
    \item d) \textbf{Sai}.
\end{itemize}
}
\end{ex}

\begin{ex} %[3D2V25]
Có 38 mét vải dùng may đồng phục, mỗi bộ may hết 3 mét vải:
\choiceTF
{\True May được nhiều nhất 12 bộ đồng phục.}
{\True Còn thừa đúng 2 mét vải.}
{Cần thêm ít nhất 2 mét vải nữa để may thêm 1 bộ nữa.}
{\True May 10 bộ đồng phục hết đúng 30 mét vải.}
\loigiai{
\begin{itemize}
    \item Ý a) Đúng: $38 : 3 = 12$ (dư 2).
    \item Ý b) Đúng: $38 - 12 \times 3 = 2\text{ m}$.
    \item Ý c) Sai: Chỉ cần thêm $3 - 2 = 1\text{ m}$ vải nữa là may thêm được 1 bộ.
    \item Ý d) Đúng: $10 \times 3 = 30\text{ m}$.
\end{itemize}
\begin{itemize}
    \item a) \textbf{Đúng}.
    \item b) \textbf{Đúng}.
    \item c) \textbf{Sai}.
    \item d) \textbf{Đúng}.
\end{itemize}
}
\end{ex}

\textbf{Phần 3. Trả lời ngắn}
\setcounterpageref{ex}{0}

\begin{ex} %[3D2V25]
Tìm số dư trong phép chia: $53 : 7$.
\shortans{4}
\loigiai{
\textbf{Đáp số:} $4$\\
$53 : 7 = 7$ (dư 4) vì $7 \times 7 + 4 = 53$.
}
\end{ex}

\begin{ex} %[3D2V25]
Trong một phép chia có số chia là 8, thương bằng 7 và số dư lớn nhất có thể có. Tìm số bị chia của phép chia đó.
\shortans{63}
\loigiai{
\textbf{Đáp số:} $63$\\
Số dư lớn nhất là 7. Số bị chia: $8 \times 7 + 7 = 56 + 7 = 63$.
}
\end{ex}

\begin{ex} %[3D2V25]
Có 43 quyển vở phát đều cho 6 bạn. Hỏi mỗi bạn được nhiều nhất bao nhiêu quyển vở?
\shortans{7}
\loigiai{
\textbf{Đáp số:} $7$\\
$43 : 6 = 7$ (dư 1). Mỗi bạn được nhiều nhất 7 quyển vở.
}
\end{ex}

\begin{ex} %[3D2V25]
Cần ít nhất bao nhiêu chiếc thuyền loại 5 chỗ ngồi để chở hết 32 du khách?
\shortans{7}
\loigiai{
\textbf{Đáp số:} $7$\\
$32 : 5 = 6$ (dư 2). Cần ít nhất $6 + 1 = 7$ chiếc thuyền.
}
\end{ex}

\begin{ex} %[3D2V25]
Tìm một số biết rằng lấy số đó chia cho 6 thì được thương là 8 và số dư là 3.
\shortans{51}
\loigiai{
\textbf{Đáp số:} $51$\\
Số đó là: $6 \times 8 + 3 = 48 + 3 = 51$.
}
\end{ex}

\begin{ex} %[3D2V25]
Một người xếp 58 quả trứng vào các vỉ trứng, mỗi vỉ để được 6 quả trứng. Hỏi sau khi xếp đầy các vỉ, còn lại bao nhiêu quả trứng chưa được vào vỉ?
\shortans{4}
\loigiai{
\textbf{Đáp số:} $4$\\
$58 : 6 = 9$ (dư 4). Còn lại 4 quả trứng.
}
\end{ex}
